\documentclass[10pt,a4paper]{article}

\usepackage[margin=1.5cm]{geometry}
\usepackage[T1]{fontenc}
\usepackage[utf8]{inputenc}
\usepackage{lmodern}
\usepackage{enumitem}
\usepackage[hidelinks]{hyperref}
\usepackage{xcolor}
\usepackage{multicol}
\usepackage{titlesec}
\usepackage{microtype}
\usepackage{parskip}
\usepackage{amsmath}

\setlength{\columnsep}{0.8cm}
\setlength{\parindent}{0pt}

\titleformat{\section}{\large\bfseries}{\thesection.}{0.5em}{}
\titleformat{\subsection}{\normalsize\bfseries}{\thesubsection.}{0.5em}{}

\setlist[itemize]{leftmargin=*, itemsep=0.2em, topsep=0.2em}
\setlist[description]{style=nextline, leftmargin=0pt, font=\normalfont\bfseries, itemsep=0.5em}

\title{\textbf{Thesis Defense --- Q\&A}}
\author{}
\date{}

\begin{document}

\maketitle

\vspace{0.5em}

\begin{multicols}{2}

\section{Detector-by-Detector Q\&A}

\begin{description}
    \item[Q: Shared Database --- How do you detect it?]
    I parse each service's \texttt{application.yml}/\texttt{.properties} for \texttt{spring.datasource.url}, resolving placeholders to defaults. Then I group services by \textbf{resolved URL} --- any URL used by more than one service is a \textbf{shared database}.

    \item[Q: Cyclic Dependency --- How do you detect it?]
    On the directed dependency graph I run \textbf{Tarjan's algorithm}; any strongly connected component with more than one service is a cycle.  \(O(V+E)\). I produce a readable chain like ``Order \(\rightarrow\) Payment \(\rightarrow\) Order'' and attach the source calls.

    \item[Q: Nano Service --- How do you detect it?]
    A service is flagged when it's below \textbf{both} thresholds: fewer than \textbf{500 LOC} \emph{and} at most \textbf{2 endpoints}. LOC comes from walking \texttt{.java} files; endpoints come from \textbf{Spoon} annotation scanning.

    \item[Q: God Service --- How do you detect it?]
    A \textbf{Spoon} analysis over six per-class metrics (\(\geq 25\) fields, \(\geq 30\) public methods, \(\geq 1000\) LOC, \(\geq 20\) import domains, \(\geq 12\) constructor params, TCC \(< 0.5\)); a class is a God Class if it exceeds at least \textbf{3 of the 6}, and a single God Class flags the service.

    \item[Q: Chatty Service --- How do you detect it?]
    Two approaches. From the dependency graph, any edge with call count \(\geq 5\), where each \texttt{@FeignClient} method counts as one call. Plus a \textbf{Spoon source scan} for chatty client types the graph misses.

    \item[Q: Hardcoded Endpoints --- How do you detect it?]
    A line-by-line scan of \texttt{.java} files for URL patterns (\texttt{http://}, \texttt{https://}, \texttt{localhost:}, \texttt{127.0.0.1}) inside string literals, extracting the full URL with a \textbf{regex}. I skip comments, imports, package declarations, and test files.

    \item[Q: Distributed Monolith --- How do you detect it?]
    A system-level check using three metrics: \textbf{coupling coefficient} \(C\), \textbf{connected ratio} \(R\), and \textbf{shared-database count} \(D\). I flag the system if $C > 0.5$ or $R > 0.8 \text{ and } D > 0$ or $R > 0.8 \text{ and } C > 0.3$. Only runs with at least 3 services.

    \item[Q: API Versioning Absence --- How do you detect it?]
    During endpoint extraction I set \texttt{hasVersioning} when any recognized strategy appears: URL path (\texttt{/v1/}), version headers, vendor media types, or query parameters like \texttt{version=}. A service is flagged when \textbf{none} of its endpoints carry a recognized version indicator.

    \item[Q: Wrong Cuts --- How do you detect it?]
    \textbf{Bidirectional dependencies}: a pair of services calling each other in both directions, which signals they are too tightly coupled to be separate. Each pair is reported once, with evidence from both directions.

    \item[Q: ESB Misuse --- How do you detect it?]
    I look for a \textbf{central mediator} using three signals against a \textbf{0.4 threshold}: high caller and callee ratios, a volume-based mediator ratio, or normalized betweenness centrality. Any one triggers a flag. \textbf{Gateway}-named services are excluded.
\end{description}

\section{Theoretical Questions}

\begin{description}
    % \item[Q: What is an anti-pattern, and which ones do you detect?]
    % An \textbf{anti-pattern} is a common solution to a recurring problem that does more harm than good. I detect ten, across four dimensions: \textbf{service design} (Nano Service, God Service), \textbf{communication} (Chatty Service, Cyclic Dependency, Hardcoded Endpoints, ESB Misuse), \textbf{data management} (Shared Database), and \textbf{deployment \& coupling} (Distributed Monolith, API Versioning Absence, Wrong Cuts).

    \item[Q: Why static analysis instead of dynamic/runtime analysis?]
    \textbf{Static analysis} needs no running system, so it works at any stage of development and integrates into \textbf{CI/CD} or code review. The tradeoff is I lose \textbf{runtime call frequencies} --- I acknowledge this in my Threats to Validity.

    % \item[Q: What makes your approach ``multi-level''?]
    % Existing tools work at one level --- either \textbf{code smells} or \textbf{architecture}. I bridge both: code-level structural metrics computed with \textbf{Spoon} --- class cohesion, field and method counts, size --- become signals for architectural-level anti-patterns like \textbf{God Service}, and \textbf{DesigniteJava}'s code-smell density feeds the \textbf{code-quality} dimension of the health score.

    \item[Q: Explain Tarjan's algorithm / how you detect cyclic dependencies.]
    I model services as a \textbf{directed graph} and run \textbf{Tarjan's algorithm} to find \textbf{strongly connected components}. Any component with more than one service is a cycle. It's a \textbf{DFS} tracking an index and a lowlink per node; runs in  \(O(V+E)\).

    \item[Q: What is the coupling coefficient?]
    It's the ratio of \textbf{actual dependency edges} to the \textbf{maximum possible}:

    $$C = \frac{|E|}{|V| \times (|V| - 1)}$$

    A high value means most services depend on most others --- a sign of a \textbf{distributed monolith}.

    \item[Q: How does ESB Misuse detection work?]
    I look for a service that mediates a disproportionate share of communication, using three signals: high \textbf{caller+callee ratios}, a volume-based \textbf{mediator ratio}, and normalized \textbf{betweenness centrality} (Brandes' algorithm). Threshold \textbf{0.4}; gateways are excluded.

    \item[Q: What is the exact research gap / novelty?]
    Existing tools usually focus on either \textbf{code-level smells} or \textbf{architecture-level dependencies}. My novelty is the \textbf{integration}: multi-level static analysis, evidence snippets, remediation, health score, and re-analysis diff in one developer-facing workflow.

    % \item[Q: Which microservice principles do the anti-patterns violate?]
    % \textbf{Shared DB}: decentralized data ownership. \textbf{Cycles, Wrong Cuts, Distributed Monolith}: independent deployability. \textbf{Chatty/ESB}: loose coupling and lightweight communication. \textbf{God/Nano}: poor service boundary design.

    \item[Q: How is this different from MSANose or MARS?]
    \textbf{MSANose} is closest, but lacks general DesigniteJava smell density, composite health score, and my God/Chatty implementation. \textbf{MARS} covers more anti-patterns, but my focus is depth: evidence snippets, scoring, re-analysis diff, and web/API workflow.

    % \item[Q: What is TCC?]
    % \textbf{Tight Class Cohesion} is the fraction of public-method pairs sharing at least one instance field. Low TCC means methods work on disjoint state, suggesting unrelated responsibilities in one class. It is one of six God Class metrics.
\end{description}

\section{Practical Questions}

\begin{description}
    % \item[Q: Walk me through what happens when a user uploads a project.]
    % The project is ingested by \textbf{ZIP upload} or \textbf{Git clone}. Then the pipeline runs asynchronously in five phases: (1) \textbf{detect microservices}, (2) run \textbf{DesigniteJava} per service, (3) build the \textbf{inter-service dependency graph} with Spoon, (4) run all \textbf{anti-pattern detectors}, (5) \textbf{assemble the result} and compute the \textbf{health score}. The frontend polls job status throughout.

    % \item[Q: How is the health score computed? Give an example.]
    % 100 points split across four capped categories: \textbf{Anti-Patterns (40), Code Quality (20), Architecture (25), Service Sizing (15)}. Each detected issue subtracts points; the score is clamped to 0--100 and mapped to a letter grade.

    % \textbf{Example --- microservice-recruit scored 40 (F):} anti-pattern penalties exceeded the 40-point budget \(\rightarrow\) 0; Code Quality 12/20; Architecture 16/25; Service Sizing 12/15.

    % $$0 + 12 + 16 + 12 = 40$$

    % \item[Q: How do you decide what is a microservice vs.\ just a module?]
    % A \textbf{deployability gate}. I scan for build files (\texttt{pom.xml}, \texttt{build.gradle}), filter out non-service folders, then require one of three signals: \texttt{@SpringBootApplication} (\textbf{HIGH}), a \textbf{Dockerfile} (\textbf{MEDIUM}), or a \texttt{main()} method (\textbf{LOW}).

    \item[Q: What does the re-analysis diff feature do?]
    It compares a new analysis to the previous one: which anti-patterns were \textbf{resolved, new, or unchanged}, plus per-metric and per-category deltas and the \textbf{health-score change}. It's inspired by \textbf{SonarQube}'s ``Clean as You Code.''

    \item[Q: How do you collect evidence for a finding?]
    Each detector extracts the relevant \textbf{source snippet} --- for a cycle, the actual \texttt{@FeignClient} or \texttt{RestTemplate} call --- with the line highlighted, plus affected services and remediation advice. It's stored as \textbf{JSON} and shown in the UI.
\end{description}

% \section{Likely ``Sneak-In'' Questions}

% \begin{description}
%     \item[Q: What's the biggest limitation?]
%     It's \textbf{Java/Spring-only}; it can't resolve \textbf{dynamically-configured URLs} (e.g.\ Spring Cloud Config); and the boundary heuristic can misclassify a shared library as a nano service.

%     \item[Q: Why was no Distributed Monolith detected?]
%     The evaluated projects showed \textbf{moderate coupling}, not the pervasive coupling required for a distributed monolith. None crossed the thresholds.

%     \item[Q: How did you validate the results?]
%     \textbf{Manual inspection} of the 98 instances --- checking datasource URLs, controller mappings, call graphs, and so on. A single reviewer (me) is a known bias.
% \end{description}

\section{Source-Code Questions}

\begin{description}
    \item[Q: Show me / explain how a detector is implemented.]
    Every detector implements the \texttt{AntiPatternDetector} interface with one method: \texttt{detect(Project, List<Microservice>)}. It's the \textbf{Strategy pattern}. Each is a Spring \texttt{@Component}, and Spring auto-collects them by injecting \texttt{List<AntiPatternDetector>} into the orchestrator. A \texttt{BaseDetector} superclass provides shared helpers.

    \item[Q: How would you add an eleventh detector?]
    Implement the interface, annotate the class \texttt{@Component}, done. Spring discovers and registers it automatically --- no orchestrator changes.

    % \item[Q: Explain the God Service detection.]
    % A \textbf{Spoon structural analysis} evaluates every class in the service against six metrics: \textbf{field count, public methods, LOC, import domains, constructor params, and TCC}. A class is a \textbf{God Class} if it exceeds at least \textbf{3 of the 6}, and a single God Class flags the service. Pure \textbf{data-holder classes} are excluded first.

    \item[Q: Per-method Feign counting.]
    Each method in a \texttt{@FeignClient} interface counts as a separate call, so a 50-method client produces a call count of \textbf{50} --- that feeds \textbf{Chatty Service} detection accurately.

    \item[Q: Async correctness.]
    I start the analysis worker via Spring's \texttt{afterCommit()} callback so the job row is committed before the background thread reads it, and progress updates run in \texttt{REQUIRES\_NEW} transactions so polling clients see them immediately.
\end{description}

\section{Additional Commission-Style Questions}

\subsection{Extra Theoretical Questions}

\begin{description}
    \item[Q: What is the difference between a code smell and an architectural anti-pattern?]
    A \textbf{code smell} is local to a class or method --- for example high complexity or low cohesion. An \textbf{architectural anti-pattern} appears at the service or system-boundary level --- for example Shared Database or Cyclic Dependency. My work connects both levels.

    \item[Q: Why are thresholds acceptable if architecture quality is context-dependent?]
    They are \textbf{heuristics}, not universal truth. They operationalize anti-pattern definitions from the literature into measurable signals. The tool detects \textbf{evidence-backed candidates}, not absolute truth.

    \item[Q: Why did you not use machine learning?]
    Mainly \textbf{explainability} and \textbf{data availability}. Publicly labeled datasets are limited, and a rule-based tool can show exactly why a service was flagged.

    \item[Q: What is the strongest threat to validity in your work?]
    \textbf{Construct validity}: whether the static signals fully capture the real architectural problem. For example, Chatty Service is approximated from static declarations, not runtime traffic.

    \item[Q: Why is API Versioning Absence considered an anti-pattern?]
    Because microservices evolve \textbf{independently}. Without explicit versioning, provider changes can break consumers and force \textbf{coordinated releases}.

    \item[Q: How did you choose the health-score category weights (40/25/20/15)?]
    \textbf{Engineering judgement}, informed by composite metrics in the literature --- SonarQube, SQALE, and \textbf{ISO 25010}. Anti-Patterns gets the largest budget because architectural issues are the focus and are expensive to fix later.

    \item[Q: Doesn't clamping a category at zero lose information?]
    Yes. That's a known limitation of the \textbf{category-cap design}. The per-issue counts are still preserved, and a density-style scaled penalty would be future work.

    \item[Q: Why does a single God Class flag the whole service?]
    The default is intentionally \textbf{sensitive}: a false positive is cheap to dismiss, while a false negative lets architectural debt grow undetected.
\end{description}

\subsection{Extra Practical Questions}

\begin{description}
    \item[Q: If a company used your tool, what should they fix first?]
    I would prioritize issues that block \textbf{independent deployment}: Cyclic Dependency, Shared Database, Wrong Cuts, and Distributed Monolith indicators.

    \item[Q: How would this integrate into CI/CD?]
    The backend already supports \textbf{asynchronous analyses} and stores history. In CI/CD, it could run after a merge request or nightly build and report \textbf{new critical findings}. A good gate is ``\textbf{no new critical issues}'' or ``health score must not regress.''

    \item[Q: What result surprised you most?]
    The frequency of \textbf{API Versioning Absence} and \textbf{Nano Service}. They show that projects can look microservice-based while still carrying architectural debt.

    \item[Q: Why did no project receive an A grade?]
    Because the score combines \textbf{multiple dimensions}. A project can avoid severe cycles and still lose points in API versioning, hardcoded endpoints, code-smell density, or service sizing.

    \item[Q: Is eleven projects enough for evaluation?]
    Enough for a \textbf{feasibility evaluation}, not a statistical claim. The corpus is heterogeneous and exercises different detectors. A true precision/recall study would need a labelled benchmark and multiple reviewers.

    \item[Q: Why did FTGO/PiggyMetrics only get C grades?]
    The score measures conformance to \textbf{this catalogue and conventions}, not overall software excellence. FTGO mainly loses points for API versioning absence; the breakdown explains whether a grade comes from a structural defect or a convention not adopted.

    \item[Q: What if API versioning is handled at a gateway or through contracts?]
    Then the detector may flag a candidate even if the team has a valid external strategy. The tool detects \textbf{service-code evidence} of versioning; it cannot prove gateway-level or consumer-contract versioning from source code alone.

    \item[Q: What would you change before deploying this in an enterprise environment?]
    I would add \textbf{organization-specific threshold profiles}, better support for \textbf{Spring Cloud Config} and \textbf{Kubernetes manifests}, and eventually \textbf{runtime telemetry}.

    \item[Q: What is the most important future improvement?]
    \textbf{Multi-language support}. Build a language-agnostic dependency graph from Docker Compose, Kubernetes, or gateway config, then add language-specific analyzers for Python, Go, .NET, Node.js, etc.

    \item[Q: How does a developer know whether a finding is a false positive?]
    Each finding includes \textbf{affected services, severity, explanation, remediation advice, and evidence snippets}. The tool gives concrete evidence, not just a black-box verdict.

    % \item[Q: How exactly does each health-score category compute its penalty?]
    % \textbf{Code Quality (20):} smell density per KLOC, scaled: $\min(20, \operatorname{round}(20 \times \text{density} / 80))$
    % \textbf{Architecture (25):} coupling coefficient above 0.1 costs
    % $\min(15, \operatorname{round}(\text{coupling} \times 15))$ plus $\min(10, 5 \text{ per cycle})$
    % \textbf{Service Sizing (15):}
    % $\min(8, 3 \text{ per nano service}) + \min(10, 5 \text{ per god service})$
    % \textbf{Anti-Patterns (40):} flat \(-8/-5/-3/-1\) by severity for everything else.

    % \item[Q: What do the HIGH/MEDIUM/LOW confidence levels actually do?]
    % Each detected service records which signal admitted it --- \textbf{framework entry point} is HIGH, \textbf{Dockerfile} MEDIUM, \texttt{main()} LOW --- and that is surfaced in the result. There are also \textbf{fallbacks} to avoid an empty analysis.

    \item[Q: What does the diff show if a service was renamed between analyses?]
    Findings are matched by a \textbf{stable key} --- anti-pattern type plus affected service plus a type-specific signature --- so a renamed service appears as \textbf{resolved + new}, not unchanged.

    % \item[Q: What does the frontend add beyond raw JSON?]
    % It gives the workflow: upload/clone, progress tracking, health score, category breakdown, expandable evidence snippets, dependency graph, JSON export, history, and re-analysis diffs.

    \item[Q: How is the application deployed?]
    Docker Compose: \textbf{Spring Boot backend}, \textbf{Angular/Nginx frontend}, and \textbf{PostgreSQL}. Production exposes only the frontend, reverse-proxies API calls internally, keeps PostgreSQL private, and externalizes secrets.
\end{description}

\subsection{Extra Source-Code Questions}

\begin{description}
    \item[Q: What are the main REST API groups?]
    \texttt{/api/auth} for register/login/current user; \texttt{/api/projects} for upload, clone, history, re-analysis and deletion; \texttt{/api/jobs} for status, results, diff, recent jobs and cancellation.

    \item[Q: How is the application secured?]
    Stateless \textbf{JWT} with Spring Security. Public endpoints: auth, Swagger, actuator health/info. Protected requests pass through a JWT filter; passwords are stored with \textbf{BCrypt}; controllers use the authenticated user.

    \item[Q: Why does each analysis job own its own microservice snapshot?]
    To preserve \textbf{history}. Re-analysis should not overwrite old services, endpoints, dependencies, smells or findings. Each \texttt{AnalysisJob} has its own snapshot, so diffs are accurate and old results remain reproducible.

    \item[Q: Why do you start the analysis worker after transaction commit?]
    Uploading or cloning creates the project and job rows in a transaction. The worker starts through an \texttt{afterCommit()} callback, so the async thread cannot read a job that has not been committed yet. That avoids a \textbf{race condition}.

    \item[Q: How is progress visible while analysis is running?]
    The worker updates the \textbf{job status} between phases, and the frontend \textbf{polls} these fields so the user sees progress instead of waiting for one blocking request.

    \item[Q: How did you protect ZIP upload extraction?]
    The extractor normalizes ZIP entry paths to prevent \textbf{Zip Slip}, rejects duplicates, streams files while counting decompressed bytes, enforces limits, and deletes the partial workspace if extraction fails.

    \item[Q: How does the dependency graph resolve a call target to a service?]
    It builds aliases from \textbf{service directory names}, \texttt{spring.application.name}, and configured ports. Then it resolves \textbf{Feign}, \textbf{RestTemplate}, and \textbf{WebClient} calls using exact, port-based, and fallback fuzzy matching.

    \item[Q: Where is API versioning detected in code?]
    \texttt{DependencyGraphBuilder} extracts endpoints and stores \texttt{hasVersioning}/\texttt{apiVersion}. \texttt{ApiVersioningDetector} then flags a service only when \texttt{versionedCount == 0}.

    \item[Q: Why use Spoon instead of regular expressions everywhere?]
    \textbf{Spoon} understands Java structure: annotations, methods, invocations, types, and class relationships. That makes it more reliable for endpoint extraction, Feign clients, TCC metrics, and God Service detection. \textbf{Regex} is only used for naturally textual signals like hardcoded URLs.

    \item[Q: How do you avoid double-counting Nano Service and God Service in the health score?]
    They are excluded from the general \textbf{Anti-Patterns} category and counted only in \textbf{Service Sizing}. That way they affect the score once, in the correct category.

    \item[Q: What happens if one detector fails?]
    The orchestrator logs that detector failure and \textbf{continues} with the remaining detectors. The benefit is \textbf{resilience}; the tradeoff is that the UI should expose partial-analysis warnings clearly.

    \item[Q: Where is the system extensible in code?]
    The main extension point is the \texttt{AntiPatternDetector} interface. A new detector implements \texttt{detect(Project, List<Microservice>)}, is annotated as a Spring component, and is automatically included.

    \item[Q: How is the re-analysis diff implemented?]
    \texttt{buildDiff()} indexes previous findings by an \textbf{issue key}: anti-pattern type + primary affected service + a type-specific signature. Each current finding looks up that key: match means \textbf{unchanged}, no match means \textbf{new}, and anything left unmatched on the previous side is \textbf{resolved}. Score deltas are recomputed through the same \texttt{HealthScoreCalculator}.

    \item[Q: Where does the deployability gate live in code?]
    \texttt{detectServicesWithConfidence()} returns \texttt{DetectedService} records with path, confidence enum, and the admitting signal. The pipeline is: scan for build files \(\rightarrow\) drop excluded folders \(\rightarrow\) drop aggregator modules \(\rightarrow\) apply the three-signal gate. Explicit fallbacks ensure a usable LOW-confidence result.

    \item[Q: Where are the score weights and thresholds defined in code?]
    Category budgets are constants in \texttt{HealthScoreCalculator}; tunable values are injected via \texttt{@Value} from \texttt{application.yml}. Detector thresholds also live in this \textbf{externalized configuration}.

    % \item[Q: How did you test the tool?]
    % Dedicated tests per detector: positive, negative and threshold-boundary cases. \texttt{BaseDetector} helpers are tested through a concrete subclass. Controller tests use \texttt{@WebMvcTest}, \texttt{MockMvc}, and mocked services.
\end{description}

\section{Additional Edge-Case Questions}

\subsection{Theoretical}

\begin{description}
    \item[Q: How sensitive are the results to default thresholds?]
    They are sensitive, especially for Nano Service, Chatty Service, God Service and ESB Misuse. The thresholds are \textbf{configurable heuristics}, not universal truth; enterprise use should calibrate them.

    \item[Q: Can one structural problem be penalized twice?]
    Yes. A bidirectional dependency can be both a \textbf{two-node cycle} and evidence of \textbf{Wrong Cuts}. I avoid double-counting Nano/God through Service Sizing, but cycle/wrong-cut overlap remains visible and should be read as an upper-bound penalty.

    \item[Q: Why use the same 0.4 ESB threshold for different signals?]
    The ESB signals are normalized to a 0--1 range, so 0.4 means a service participates in a large share of communication. It is a pragmatic default; small graphs can produce candidate mediators more easily.

    \item[Q: Could God Service detection miss a broad service with no God Class?]
    Yes. The detector catches responsibility concentration visible in one oversized/low-cohesion class. A service with many medium-sized but unrelated classes would need service-level cohesion or domain analysis.

    \item[Q: What does the coupling coefficient hide?]
    It compresses the graph to density, so it loses dependency type, business criticality and runtime frequency. It is one signal, not a full architectural judgment.

    \item[Q: Can DesigniteJava overestimate smell density?]
    Yes. Since each service is analyzed separately, cross-service DTOs or abstractions may look unused. I treat smell density as a \textbf{proxy} and possible upper bound.

    \item[Q: What claims does the evaluation support?]
    It supports \textbf{feasibility}: the tool analyzes real projects and produces differentiated, evidence-backed findings. It does not prove statistical precision/recall or industry-wide prevalence.

    \item[Q: Why include a polyglot repository?]
    Real microservice repositories are often mixed-language. The tool analyzes the Java-visible part and exposes the limitation that non-Java services are outside scope.

    \item[Q: Does gateway exclusion hide real ESB bottlenecks?]
    It can. Gateways are excluded because centrality is expected there; otherwise valid gateway architectures would be flagged. A production profile could make this configurable.

    \item[Q: How do you compare with MARS if MARS detects more patterns?]
    MARS has broader catalogue coverage. My contribution is the \textbf{integrated workflow}: boundary detection, evidence, scoring, diffing and a web/API application.
\end{description}

\subsection{Practical}

\begin{description}
    \item[Q: Were evaluated repository versions pinned?]
    For strict reproducibility they should be pinned to commit SHAs. If only default branches were cloned, results are a snapshot at analysis time and may change as repositories evolve.

    \item[Q: What were the dataset selection criteria?]
    Public GitHub/GitLab availability, Java/JVM microservice code, Maven or Gradle build files, and at least three services. I also chose varied sizes and domains to exercise different detectors.

    \item[Q: How would enterprise results differ?]
    Enterprise systems may have more services, historical debt, messaging, external config and polyglot code. That could reveal more coupling but also more false negatives for the current static Java/Spring scope.

    \item[Q: What about Spring profiles or external config for database URLs?]
    The detector resolves simple placeholders/defaults, but not full runtime profile activation or Spring Cloud Config. Shared Database detection can miss values resolved only at deployment time.

    \item[Q: What if DesigniteJava is missing or times out?]
    The analysis continues. Missing JAR means skipped analysis; timeout stops the subprocess; non-zero exit still allows parsing any produced CSV. Code Quality may then have fewer smell records.

    \item[Q: Is cancellation immediate?]
    It is \textbf{cooperative}. The worker checks cancellation before phases and between per-service steps, but it does not instantly interrupt every low-level Spoon or Designite operation.

    \item[Q: How do you prevent cross-user access by ID?]
    Controllers use the authenticated principal, and services query/check resources through the owning user or project. An ID alone is not enough to access another user's project or job.

    \item[Q: Why only public GitHub/GitLab HTTPS repos?]
    It keeps ingestion safe and simple for the prototype. Private repos require token handling, secret storage, revocation and audit controls.

    \item[Q: What are the JWT browser-storage tradeoffs?]
    It is simple for an SPA, but local storage is exposed if XSS exists. Production could use shorter-lived tokens, CSP, refresh-token rotation or secure HttpOnly cookies.

    \item[Q: Why disable CSRF and CORS?]
    CSRF is disabled because the API uses stateless bearer tokens, not cookie sessions. CORS is disabled because production is expected to route browser traffic through the frontend/proxy; separate origins need explicit CORS config.

    \item[Q: Are diffs comparable if thresholds change?]
    Only partly. The diff compares two configured runs, not a pure source-code delta. Strict comparison requires keeping or versioning the same threshold profile.

    \item[Q: What CI gate would you use?]
    Start with \textbf{no new critical findings} and \textbf{score must not regress}, not failing on every legacy issue.
\end{description}

\subsection{Source Code}

\begin{description}
    \item[Q: Where is user ownership enforced?]
    In the service layer, especially \texttt{ProjectService} and \texttt{JobService}: resources are loaded by owner or checked through the owning project before returning data.

    \item[Q: How is cancellation implemented in code?]
    \texttt{JobService.cancelJob()} marks the job \texttt{CANCELLED}. \texttt{AnalysisWorker} calls \texttt{isCancelled(jobId)} before and between major phases and stops when it sees that status.

    \item[Q: What if Spoon fails to parse a service/class?]
    The scanner logs the error and skips that model or class, favouring resilience. The risk is false negatives, so production should surface parse warnings.

    \item[Q: How does fuzzy dependency matching avoid false positives?]
    It tries safer strategies first: aliases, \texttt{spring.application.name}, and port matching. Fuzzy matching is a fallback and should be interpreted with the stored evidence.

    \item[Q: Why store evidence/details as JSON text?]
    Detector evidence has different shapes. JSON keeps detector-specific payloads flexible while common fields remain relational. The tradeoff is weaker database-level structure for details.

    \item[Q: How do disabled detectors affect the score?]
    Disabled detectors produce no findings, so they create no penalties. The score reflects the active detector profile, not an absolute full-catalogue assessment.

    \item[Q: What happens under load?]
    Analyses run through a configured bounded async executor. Extra jobs wait in the queue; production should add rate limits and worker scaling.

    \item[Q: Why configure pool and queue sizes?]
    Static analysis is CPU/I/O-heavy. A bounded executor prevents uploads from exhausting the server and lets deployment match hardware capacity.

    \item[Q: Can the diff key misclassify findings?]
    Yes, in edge cases. The key uses type, affected services and signatures, and duplicates are matched one-to-one with a deque, but renames or weak signatures can still confuse matching.

    \item[Q: Is persisted score authoritative?]
    The score is stored, but breakdowns/diffs are recomputed by \texttt{HealthScoreCalculator}. If scoring logic changes, old breakdowns can differ unless scoring profiles are versioned.

    \item[Q: Why count LOC from \texttt{.java} files only?]
    Java LOC matches the Java/Spring analysis scope. It excludes YAML, SQL, frontend code and non-Java services, so it is incomplete for polyglot systems.

    \item[Q: How would you test a new detector?]
    Positive, negative, threshold-boundary, false-positive, malformed-input and evidence-snippet tests; for graph detectors, also small/disconnected/duplicate-edge graphs.
\end{description}

\section{Challenge Angles}

\begin{description}
    \item[Q: Why should developers trust the health score?]
    They should not treat it as absolute truth. It is a \textbf{summary indicator}; the real value is the transparent breakdown, where every deduction is traceable to a detector, severity, affected services and evidence.

    \item[Q: What is your weakest assumption?]
    That static source-level signals are good proxies for runtime behaviour. Static Feign method count approximates chatty communication, but real traffic volume needs runtime telemetry.

    \item[Q: Why source code instead of bytecode?]
    Source code gives access to \textbf{config files, annotations, endpoint mappings, snippets and remediation evidence}. Bytecode is useful for dependencies, but weaker for developer-facing explanations.

    \item[Q: What about Kafka/RabbitMQ/asynchronous messaging?]
    Current graph extraction targets REST-style calls: Feign, RestTemplate and WebClient. Messaging dependencies are a limitation and a natural future extension.

    % \item[Q: Why PostgreSQL?]
    % The model is relational: users, projects, jobs, services, endpoints, dependencies, smells, results and findings. PostgreSQL fits historical analysis and diff persistence well.

    \item[Q: Why persist snippets?]
    Workspaces may be deleted or changed. Persisting snippets during analysis makes results reproducible even if the source checkout is gone.

    \item[Q: False positives vs.\ false negatives?]
    False positives: thresholds and context-sensitive detectors; evidence lets users dismiss them quickly. False negatives: dynamic URLs, external config, runtime-only traffic, messaging and non-Java services.

    \item[Q: Most context-sensitive vs.\ objective detectors?]
    \textbf{Context-sensitive}: Nano Service and API Versioning Absence. \textbf{More objective}: Shared Database and Cyclic Dependency.

    % \item[Q: Is this just a wrapper around existing tools?]
    % No. \textbf{DesigniteJava} is one input. The contribution is the full pipeline: boundary detection, Spoon endpoint/dependency extraction, graph detectors, custom detectors, evidence, scoring, persistence, diffing and UI workflow.

    % \item[Q: Isn't the health score subjective?]
    % Yes, like most composite quality scores. The value is that it is \textbf{transparent, decomposed, repeatable}, and useful for comparing runs over time.

    % \item[Q: If the demo fails?]
    % Say: ``The live run appears slow, so I'll switch to a completed analysis of the same flow.'' Show the backup result; do not debug live.

    % \item[Unexpected wording structure]
    % Use: \textbf{input \(\rightarrow\) method \(\rightarrow\) tradeoff \(\rightarrow\) limitation}. Example: dependency graph \(\rightarrow\) Tarjan SCC \(\rightarrow\) static calls not runtime traffic \(\rightarrow\) dynamic/message calls can be missed.
\end{description}

% \section{If I Blank}

% \begin{description}
%     \item[Q: What can I say if I blank?]
%     \begin{itemize}
%         \item ``Could you repeat the question?''
%         \item ``Let me start from the high level: \ldots''
%         \item ``The exact value I'd need to check, but the idea is \ldots''
%         \item ``That's in Chapter X --- the short version is \ldots''
%     \end{itemize}
% \end{description}

\end{multicols}

\end{document}
